\textbf{Время:} 01:00

\textbf{Тезис:} сравнение должно относиться к одинаково определённым наблюдаемым и классам пар.

\textbf{Визуал:} сопоставить строки таблицы эксперимента и модели.

\textbf{Текст:} Одинаковых Au+Au и трёх ГэВ недостаточно. Нужно согласовать заряды, интервалы импульса, центральность и продольные отборы. Быстрота пары и псевдобыстрота отдельной частицы — разные ограничения.

Далее выбираем одинаковую систему отсчёта и размерность анализа. Одномерный R-inv нельзя напрямую считать трёхмерным радиусом STAR. Проверяем взаимодействия: у показанного эксперимента есть учёт Кулона, простой косинус модели его не включает.

Наконец, нужны проверенные ошибки, нормировка и отбор событий. Тогда различие сможет ограничить физическую модель, а не просто показать, что две программы посчитали разные величины.

\textbf{Переход:} Вернёмся к исходному вопросу о размерах источника, который уже исчез.
